% Faithful formal-publication transcription: group targets and complete constructive Baer context/continuations, pp16:3 and16:29--31.
\section*{Group-isomorphism targets}
\begin{authorclaim}{TheoremA}
Let $p$ be an odd prime, and let $\operatorname{GpIso2Exp}(p)$ denote the isomorphism problem for $p$-groups of class 2 and exponent $p$ in the model of matrix groups over $\mathbb F_p$. For groups of order $p^n$, there is a search-to-decision reduction for $\operatorname{GpIso2Exp}(p)$ running in time $p^{O(n)}=\operatorname{poly}(|G|)$.
\end{authorclaim}
\begin{authorclaim}{TheoremB}
For $p$ an odd prime, $p\ge n^{\Omega(1)}$, there is a randomized counting-to-decision reduction for $\operatorname{GpIso2Exp}(p)$ for groups of order $p^n$, running in time $p^{O(n)}=\operatorname{poly}(|G|)$.
\end{authorclaim}
\section*{6.2 Constructive Baer Correspondence and Theorems A and B}
Let us review Baer's Correspondence [7], which connects alternating bilinear maps with $p$-groups of class 2 and exponent $p$. Let $P$ be a $p$-group of class 2 and exponent $p$, $p>2$. Suppose the commutator subgroup $[P,P]\cong\mathbb Z_p^m$ and $P/[P,P]\cong\mathbb Z_p^n$. Then the commutator map $[,]:P/[P,P]\times P/[P,P]\to[P,P]$ is an alternating bilinear map. Conversely, let $\phi:\mathbb Z_p^n\times\mathbb Z_p^n\to\mathbb Z_p^m$ be an alternating bilinear map. Then a $p$-group of class 2 and exponent $p$, denoted as $P_\phi$ can be defined as follows. The group elements are from $\mathbb Z_p^n\times\mathbb Z_p^m$, and the group product $\cdot$ is defined as
\[(u,v)\cdot(u',v')=(u+u',v+v'+\tfrac12\phi(u,u')).\]
We say that $(A,B)\in\mathrm{GL}(n,p)\times\mathrm{GL}(m,p)$ is a pseudo-autometry of $\phi$, if $\phi=B\circ\phi\circ A$. Wilson [71] elucidated the structure of $\operatorname{Aut}(P_\phi)$ in terms of the pseudo-autometry group of $\phi$, that we denote $\Psi\operatorname{Aut}(\phi)$. Here we recall the consequence of Wilson's result that we need for counting group isomorphisms.
\begin{authorclaim}{Proposition23 (Wilson [71, Prop.3.8], see [15, Prop.2.4] for notation closer to ours)}
For $\phi:\mathbb Z_p^n\times\mathbb Z_p^n\to\mathbb Z_p^m$ an alternating bilinear map,
\[|\operatorname{Aut}(P_\phi)|=|\Psi\operatorname{Aut}(\phi)|p^{nm},\]
where $\Psi\operatorname{Aut}(\phi)$ denotes the pseudo-autometry group of $\phi$.
\end{authorclaim}
We then state a lemma which can be viewed as a constructive version of Baer's Correspondence, communicated to us by James B. Wilson.
\begin{authorclaim}{Lemma24 (Constructive version of Baer's Correspondence for matrix groups)}
Let $p$ be an odd prime. Over the finite field $\mathbb F=\mathbb F_{p^e}$, Alternating Matrix Space Isometry is equivalent to Group Isomorphism for matrix groups over $\mathbb F$ that are $p$-groups of class 2 and exponent $p$. More precisely, there are functions computable in time $\operatorname{poly}(n,m,\log|\mathbb F|)$:
\begin{itemize}
\item $G:\Lambda(n,\mathbb F)^m\to\mathrm M(n+m+1,\mathbb F)^{n+m}$ and
\item $\operatorname{Alt}:\mathrm M(n,\mathbb F)^m\to\Lambda(m,\mathbb F)^{O(m^2)}$
\end{itemize}
such that: (1) for an alternating bilinear map $\mathbf A$, the group generated by $G(\mathbf A)$ is the Baer group corresponding to $\mathbf A$, (2) $G$ and Alt are mutually inverse, in the sense that the group generated by $G(\operatorname{Alt}(M_1,\ldots,M_m))$ is isomorphic to the group generated by $M_1,\ldots,M_m$, and conversely $\operatorname{Alt}(G(\mathbf A))$ is pseudo-isometric to $\mathbf A$.
\end{authorclaim}
\begin{proof}
First, let $G$ be a $p$-group of class 2 and exponent $p$ given by $m$ generating matrices of size $n\times n$ over $\mathbb F$. Then from the generating matrices of $G$, we first compute a generating set of $[G,G]$, by just computing all the commutators of the given generators. We can then remove those redundant elements from this generating set in time $\operatorname{poly}(\log|[G,G]|,\log|\mathbb F|)$, using Luks' result on computing with solvable matrix groups [51]. We then compute a set of representatives of a non-redundant generating set of $G/[G,G]$, again using Luks's aforementioned result. From these data we can compute an alternating bilinear map representing the commutator map of $G$ in time $\operatorname{poly}(n,m,\log|\mathbb F|)$.
Conversely, let an alternating bilinear map be given by $\mathbf A=(A_1,\ldots,A_m)\in\Lambda(n,\mathbb F)^m$. From $\mathbf A$, for $i\in[n]$, construct $B_i=[A_1e_i,\ldots,A_me_i]\in\mathrm M(n\times m,\mathbb F)$, where $e_i$ is the $i$th standard basis vector of $\mathbb F^n$. That is, the $j$th column of $B_i$ is the $i$th column of $A_j$. Then for $i\in[n]$, construct
\[\widetilde B_i=\begin{bmatrix}1&e_i^t&0\\0&I_n&B_i\\0&0&I_m\end{bmatrix}\in\mathrm{GL}(1+n+m,\mathbb F),\]
where $e_i\in\mathbb F^n$, and for $j\in[m]$, construct
\[\widetilde C_j=\begin{bmatrix}1&0&e_j^t\\0&I_n&0\\0&0&I_m\end{bmatrix}\in\mathrm{GL}(1+n+m,\mathbb F),\]
where $e_j\in\mathbb F^m$. Let $G(\mathbf A)$ be the matrix group generated by $\widetilde B_i$ and $\widetilde C_j$. Then it can be verified easily that, $G(\mathbf A)$ is isomorphic to the Baer group corresponding to the alternating bilinear map defined by $\mathbf A$. In particular, $[G,G]\cong\mathbb F^m\cong\mathbb Z_p^{em}$ (isomorphism of abelian groups), and $G/[G,G]\cong\mathbb F^n\cong\mathbb Z_p^{en}$. This construction can be done in time $\operatorname{poly}(n,m,\log|\mathbb F|)$.
\end{proof}
\begin{proof}[Proof of TheoremA]
The search-to-decision reduction follows from TheoremA', using the $q^{O(n+m)}$-time algorithm, with the constructive version of Baer's Correspondence in the model of matrix groups over finite fields (Lemma24).
In more detail, given Lemma24 we can follow the procedure in the proof of TheoremA'. For the given $p$-groups, we compute their commutator maps. Then whenever we need to feed the decision oracle, we transform from the alternating bilinear map to a generating set of a $p$-group of class 2 and exponent $p$ with this bilinear map as the commutator map. After getting the desired pseudo-isometry for the alternating bilinear maps, we can easily recover an isomorphism between the originally given $p$-groups.
\end{proof}
\begin{proof}[Proof of TheoremB]
For the counting-to-decision reduction, we basically follow the above routine, but with a twist, because of the minor distinction between alternating matrix space isometry, and alternating bilinear map pseudo-isometry. Let us briefly explain this issue. Suppose from an alternating bilinear map $\phi:\mathbb Z_p^n\times\mathbb Z_p^n\to\mathbb Z_p^m$ we constructed the $p$-group $P_\phi$ of class 2 and exponent $p$; by Proposition23 $|\operatorname{Aut}(P_\phi)|=p^{nm}|\Psi\operatorname{Aut}(\phi)|$, so by multiplying the result by $p^{nm}$, it is necessary and sufficient to count the psuedo-autometries of $\phi$.
Towards that end, let $(C_1,\ldots,C_m)\in\Lambda(n,p)$ be a matrix representation of $\phi$. If $C_i$'s are linearly independent, then for a pseudo-autometry $(A,B)\in\mathrm{GL}(n,p)\times\mathrm{GL}(m,p)$, given $A$ there exists a unique $B$ that makes $(A,B)$ a pseudo-autometry. If $C_i$'s are not linearly independent, say the linear span of $C_i$'s is of dimension $m'$, then the number of $B$ such that $(A,B)$ is a pseudo-autometry (assuming there are any) is $|\mathrm M((m-m')\times m',p)||\mathrm{GL}(m-m',p)|=p^{m'(m-m')}|\mathrm{GL}(m-m',p)|$. To see this, suppose that we have taken linear combinations of the $C_i$ so that $C_1,\ldots,C_{m'}$ are linearly independent and $C_{m'+1},C_{m'+2},\ldots,C_m$ are zero. Then without changing the $C_i$, we may take any invertible linear combination among $C_{m'+1},\ldots,C_m$ (a copy of $\mathrm{GL}(m-m',p)$), and we may add any linear combination of the last $m-m'$ matrices to the first $m'$ matrices (a copy of $\mathrm M((m-m')\times m',p)$). The counting to decision reduction for Alternating Matrix Space Isometry computes the number of $A\in\mathrm{GL}(n,p)$ so that there exists some $B\in\mathrm{GL}(m,p)$ such that $(A,B)$ is a pseudo-autometry. So it needs to be multiplied by a factor of $p^{m'(m-m')}|\mathrm{GL}(m-m',p)|$.
\end{proof}
